\documentclass[12pt,a4paper]{article}
\usepackage[utf8]{inputenc}
\usepackage[T1]{fontenc}
\usepackage[brazil]{babel}
\usepackage{lmodern}
\usepackage{amsmath}
\usepackage{geometry}
\usepackage{hyperref}
\geometry{margin=2.5cm}
\title{Relatório de conclusões --- escalonamento dinâmico de tarefas}
\author{MVHSN}
\date{06 de agosto de 2026}
\begin{document}
\maketitle

\section{Resultado do estudo}
A equação original é matematicamente consistente como modelo de minimização do \emph{makespan}. Ela é superior a uma ETA baseada apenas em médias porque considera tarefas pendentes, tarefas em execução, velocidades diferentes e a possibilidade de aguardar.

\section{Limitação identificada}
A solução exata exigiria reotimizar uma alocação heterogênea de centenas de tarefas a cada liberação. Isso é caro e pouco necessário para a operação. Além disso, os custos das tarefas não são determinísticos.

\section{Fórmula recomendada}
Cada worker é tratado como um slot independente. Para a tarefa $i$ no worker $s$, o tempo previsto é:
\[
\widehat{p}_{i,s}=\frac{\widehat{C}(k_i)\,m_i}{v_s}.
\]
Aqui, $\widehat{C}(k_i)$ é o custo aprendido do cenário, $m_i$ representa a complexidade observada da tarefa e $v_s$ é a velocidade histórica do worker.

Quando um worker fica livre, o coordenador avalia uma janela de candidatos. Para cada candidato, simula uma distribuição gulosa das tarefas restantes e calcula:
\[
\widehat{T}_{i\rightarrow s}=\max_r L'_r.
\]
A tarefa escolhida é a que minimiza $\widehat{T}_{i\rightarrow s}$. A máquina só permanece ociosa quando esperar produz ganho maior que a margem de erro da previsão.

\section{Melhorias incorporadas}
O modelo recomendado:
\begin{itemize}
\item diferencia $n=50$, $n=100$, $n=200$ e $n=400$;
\item aprende com os tempos reais;
\item trata workers individualmente;
\item considera iterações ECM e fallbacks;
\item limita a busca a candidatos relevantes;
\item usa margem de segurança;
\item mantém fallback para a política atual quando há poucos dados.
\end{itemize}

\section{Exclusão deliberada}
Temperatura, clock e carga térmica não fazem parte desta versão, conforme decisão do projeto. A velocidade será aprendida exclusivamente do histórico de execução e da capacidade observada dos workers.

\section{Validação necessária}
Ainda não é possível afirmar um percentual de ganho sem backtest. A comparação deve medir makespan, erro da ETA, tempo ocioso, vazão por cenário, equilíbrio entre $n=50$ e $n=400$, reexecuções e leases expirados.

\section{Parecer final}
A proposta deve ser aprovada como base de evolução, mas implementada como heurística adaptativa por worker, cenário e custo aprendido, com horizonte móvel, margem de segurança e fallback. Essa versão preserva a robustez da fila e permite medir o ganho antes de alterar o coordenador em produção.
\end{document}
